\rulesection{1.0 Introduction}

\textit{Shenandoah: Valley of Fire} simulates the 1862 and 1864 campaigns for the Shenandoah Valley. The 1862 campaign featured the brilliant maneuvers of Thomas Jackson against a Union force roughly double his force's size. In 1864, Phil Sheridan would enter the fray and permanently remove the valley from the Confederacy's plans. The game shows both sides' command problems. The system is based on command points and leader command ratings, which together determine movement allowances.

If you have never played a wargame before, the rules may seem long and complicated, but they merely cover simple concepts that should become clear after a few readings. The rules are not complex, but it is important that players read all of them before attempting to play.

Each section of the rules is numbered, and important paragraphs within a section carry a second number, like \texttt{2.2}. When a section contains subsections, these are identified like \texttt{2.24}. When a rule refers to another related rule, the number of that rule appears in parentheses.
